% TOP_PROBLEM: {"id": 3220, "title": "A homomorphism problem for flows", "category_id": 3, "category": {"id": 3, "name": "graph_theory", "display_name": "Graph Theory", "description": "Problems involving graphs, networks, and their properties.", "slug": "graph-theory", "order_index": 3, "created_at": "2026-07-31T15:26:25.671Z"}, "status": "open", "problem_number": "OPG-134", "set_id": 12, "statusQualification": "Makes the standard zero-free connection-set convention explicit; allowing zero while declaring Cayley graphs simple creates a degenerate false variant."}
% FIRST_POSTED: 2026-10-01
% ENTRY_KIND: statement_only
% UnsolvedMath ID 3220; source catalogue code OPG-134
% !TeX program = lualatex
% Definition and sources only; this file is not an LLM solution attempt.
\documentclass[11pt,a4paper]{article}
\usepackage[margin=25mm]{geometry}
\usepackage{fontspec}
\setmainfont{lmroman10-regular.otf}[BoldFont=lmroman10-bold.otf,
 ItalicFont=lmroman10-italic.otf,BoldItalicFont=lmroman10-bolditalic.otf]
\usepackage{amsmath,amssymb,mathtools}
\usepackage{unicode-math}
\setmathfont{latinmodern-math.otf}
\AtBeginDocument{\let\setminus\smallsetminus}
\usepackage{xurl,hyperref}
\hypersetup{colorlinks=true,urlcolor=blue}
\setcounter{secnumdepth}{0}
\setlength{\parindent}{0pt}
\setlength{\parskip}{5pt}
\setlength{\emergencystretch}{3em}
\begin{document}
\section*{A homomorphism problem for flows}\label{3220}
{\small UnsolvedMath ID 3220 · OPG-134 · Graph Theory · OpenGarden}
\subsection{Definitions and mathematical statement}
Let $M$ and $M'$ be abelian groups written additively, and let $B\subseteq M\setminus\{0\}$ and $B'\subseteq M'\setminus\{0\}$ satisfy $B=-B$ and $B'=-B'$. Their Cayley graphs have vertex sets $M$ and $M'$, respectively, with distinct $x,y$ adjacent when $y-x$ belongs to the indicated set. A graph homomorphism $h:\operatorname{Cay}(M,B)\to\operatorname{Cay}(M',B')$ is an arbitrary vertex map preserving adjacency; it need not be additive or injective.

For a finite graph $G$ with an arbitrary orientation, a $B$-flow is a map $f:E(G)\to B$ such that at every vertex the sum of values on outgoing edges equals the sum on incoming edges, with sums evaluated in $M$. Define a $B'$-flow similarly using $M'$. Symmetry of the allowed sets makes existence independent of the chosen orientation: reversing an edge changes its value to its negative.

The conjecture asks whether existence of a homomorphism between the two Cayley graphs implies the following implication for every finite $G$:
\[
       G\text{ has a }B\text{-flow}
       \ \Longrightarrow\ G\text{ has a }B'\text{-flow}.
\]
The prospective new flow need not be the edgewise image under $h$, since $h$ is not assumed to preserve group addition. The zero-free convention gives the intended nowhere-zero flow problem and keeps the Cayley graphs loopless; otherwise a zero in an allowed set would require a separate loop convention.
\subsection{Short English statement}
Does a homomorphism between two abelian Cayley graphs guarantee that every flow using the first allowed value set can be replaced by one using the second?
\subsection{Sources}
\begin{itemize}
\item[S1] ULAM.AI and the UnsolvedMath Contributors, supplied problems.json, record 3220 (OPG-134), OpenGarden. Catalogue identity and source transcription; CC BY 4.0.
\url{https://huggingface.co/datasets/ulamai/UnsolvedMath}.
\item[S2] Open Problem Garden, A homomorphism problem for flows. Original problem and discussion.
\url{http://www.openproblemgarden.org/op/a_homomorphism_problem_for_flows}.
\end{itemize}
\end{document}
